\chapter{Admissibility and the Space of Continuations}
\label{ch:admissibility}

\begin{quotation}
\noindent\textit{Synopsis.} Having defined what can be lost, this chapter
turns to what remains possible. Using concepts imported from Admissibility
Theory, it defines admissible continuation sets and task-relative
sufficiency: a substrate is never sufficient in the abstract, only relative
to a particular task. The chapter argues that sufficiency cannot generally
be certified at creation time, since future tasks are often unknown then,
and introduces the distinction between closed hidden-state domains, where
omitted information can in principle be enumerated and restored, and open
domains, where no finite augmentation can exhaust future informational
needs. This distinction becomes one of the framework's central organizing
principles.
\end{quotation}

\section{Admissibility Theory, Imported and Specialized}

Admissibility Theory supplies the second basic object this monograph
needs: given a substrate, the set of continuations still reachable from
it. As with Distinguishability Geometry in Chapter~\ref{ch:distinctions},
nothing in the prior formalism is altered here; it is specialized to the
continuation problem of \cref{ch:continuation-problem}.

\begin{definition}[Admissible Continuation Set]
\label{def:admissible-set}
$\admis{\tau}{r}$, the admissible continuation set of task $\tau$ given
rendering $r$, is the set of continuations of $\tau$ that remain reachable
given only $r$.
\end{definition}

This definition leaves the internal structure of a ``continuation''
deliberately unspecified, and deliberately so: a continuation is a legal
move in Chapter~\ref{ch:chess}'s setting, a plausible causal narrative in
Chapter~\ref{ch:painting}'s, a verified-or-rejected verdict in
Chapter~\ref{ch:proofs}'s, and a binding ruling in
Chapter~\ref{ch:contracts}'s. The formalism is domain-agnostic on purpose;
each worked domain in Part~\ref{part:worked-domains} supplies its own
concrete continuation space rather than inheriting one from this chapter.

\begin{proposition}[Admissibility Bounded by Full Access]
\label{prop:admissibility-upper-bound}
For any rendering $r$ of an object $e$ and any task $\tau$,
$\admis{\tau}{r} \subseteq \admis{\tau}{e}$.
\end{proposition}

\begin{proofsketch}
A rendering (\cref{def:rendering}) can only identify distinctions present
in the source, never introduce new ones; any continuation reachable using
only $r$ is therefore reachable using the full source $e$ as well, since
$e$ determines everything $r$ does and possibly more. The reverse
containment need not hold, which is exactly what makes discard
consequential.
\end{proofsketch}

\begin{proposition}[Admissibility Is Monotone Under Further Rendering]
\label{prop:admissibility-monotone}
If $r'$ is itself a rendering of $r$ -- that is, $r'$ is obtained from $r$
by a further projection in the sense of \cref{def:rendering} -- then
$\admis{\tau}{r'} \subseteq \admis{\tau}{r}$.
\end{proposition}

\begin{proofsketch}
Immediate from \cref{prop:admissibility-upper-bound} applied with $r$
playing the role of the source and $r'$ the rendering of it: $r'$ can only
identify distinctions already present or absent in $r$, and cannot
recover any continuation $r$ itself could not reach.
\end{proofsketch}

\cref{prop:admissibility-monotone} is the admissibility-theoretic
counterpart of Chapter~\ref{ch:budget-relay}'s
\cref{prop:monotonicity}: continuation possibilities can only shrink or
stay level as a substrate is further rendered, never expand, in the
absence of a preserved earlier stage to re-derive from
(\cref{prop:root-recovery}).

\section{What a Task Requires}

Chapter~\ref{ch:distinctions} spoke of a distinction being ``required'' by
a task $\tau$ without saying precisely what that means. The remainder of
this chapter makes that precise, and in doing so corrects an
oversimplification present in an earlier formulation of task-relative
sufficiency: that formulation spoke of ``the true or intended
continuation'' of a task, in the singular, which is adequate for tasks
like verification that admit exactly one correct outcome but is not
adequate for tasks like Chapter~\ref{ch:painting}'s evocative
reconstruction, which by design tolerates more than one acceptable
completion.

\begin{definition}[Correctness Set]
\label{def:correctness-set}
A task $\tau$, evaluated against an object $e$, determines a nonempty
correctness set $\mathcal{C}_\tau(e)$: the set of continuations that count
as successfully performing $\tau$. $\mathcal{C}_\tau(e)$ is a singleton for
tasks demanding a unique correct outcome (verification against a fixed
proof, say) and may contain many continuations for tasks tolerating
multiple acceptable completions (evocative reconstruction from a
painting, say, where several causal histories are equally acceptable).
\end{definition}

\begin{definition}[Task-Relative Sufficiency, Revised]
\label{def:sufficiency}
$r$ is sufficient for $\tau$ (evaluated against the true object $e$) iff
\[
\admis{\tau}{r} \;\subseteq\; \mathcal{C}_\tau(e).
\]
\end{definition}

This revises, and is intended to supersede, the earlier singular-outcome
phrasing. The subset formulation in \cref{def:sufficiency} is chosen
because it correctly unifies the two cases an adequate definition needs to
cover simultaneously. For a verification-type task, $\mathcal{C}_\tau(e)$
is a singleton, and \cref{def:sufficiency} demands that $r$ pin down
exactly that outcome: $\admis{\tau}{r}$ must not merely intersect the
correct answer but be contained in it, ruling out any admissible but wrong
continuation being reachable alongside the right one. For a task
tolerating multiple acceptable completions, $\mathcal{C}_\tau(e)$ is
larger, and \cref{def:sufficiency} only requires that every continuation
$r$ makes reachable falls inside that larger acceptable set -- which is
satisfied even when several distinct, mutually incompatible continuations
remain open, exactly the situation Chapter~\ref{ch:painting} calls benign
underdetermination.

\begin{proposition}[Benign Loss Preserves Sufficiency]
\label{prop:benign-preserves-sufficiency}
If every distinction in $\discard{r}$ is benign relative to $\tau$
(\cref{def:benign-loss}), and $e$ itself would be sufficient for $\tau$,
then $r$ is sufficient for $\tau$.
\end{proposition}

\begin{proofsketch}
By \cref{def:benign-loss}, no distinction required to determine membership
in $\mathcal{C}_\tau(e)$ is among those $r$ discards; hence every
continuation in $\admis{\tau}{r}$ is also determinable from $e$ as
belonging to $\mathcal{C}_\tau(e)$, since the discarded distinctions play
no role in that determination. $\admis{\tau}{r} \subseteq
\admis{\tau}{e} \cap (\text{continuations consistent with } \tau)
\subseteq \mathcal{C}_\tau(e)$ whenever $e$ is itself sufficient, giving
\cref{def:sufficiency} directly.
\end{proofsketch}

\cref{prop:benign-preserves-sufficiency} is the formal statement of the
coherence this monograph's two central chapters are supposed to have with
each other: Chapter~\ref{ch:distinctions}'s loss-level vocabulary
(benign versus corrupting) and this chapter's continuation-level
vocabulary (admissible versus sufficient) are not two independent
frameworks that happen to agree on examples; benign loss is, by
construction, exactly the kind of loss that cannot violate sufficiency.

\section{Why Sufficiency Cannot Be Verified at Production Time}

\begin{proposition}[Task Unavailability at Ingestion]
\label{prop:task-unavailability}
If the future task space $T$ relevant to a substrate is open -- unbounded,
not fixed at the time the substrate is produced -- then sufficiency of a
rendering $r$ relative to $T$ cannot be certified when $r$ is produced.
\end{proposition}

\begin{proofsketch}
Suppose otherwise: suppose $r$ could be certified sufficient for every
$\tau \in T$ at production time. Then $\admis{\tau}{r} \subseteq
\mathcal{C}_\tau(e)$ for every $\tau \in T$, which requires $r$ to have
preserved every distinction any such $\tau$ might need to rule out an
incorrect continuation. Since $T$ is open, a task $\tau^*$ can always be
constructed, after the fact, requiring a distinction not represented in
$r$ -- for instance, a distinction $r$'s producer had no reason to
anticipate -- for which $\admis{\tau^*}{r} \not\subseteq
\mathcal{C}_{\tau^*}(e)$. This contradicts the certification assumed
above. Hence no such certification is available in general.
\end{proofsketch}

This result is the formal version of \cref{principle:global-insufficiency}
for a single substrate. It motivates the shift, completed in
\cref{def:root-object}, from optimizing a rendering for sufficiency --
which \cref{prop:task-unavailability} shows to be generally impossible to
verify -- to optimizing it for re-derivability from a preserved root
object.

\section{Closed and Open Hidden-State Domains}

\begin{definition}[Closed versus Open Hidden-State Domain]
\label{def:closed-open}
A domain has \emph{closed} hidden state if the distinctions omitted from a
given state are drawn from a finite, enumerable list -- for instance, a
chess position's castling rights, en passant eligibility, and repetition
count. A domain is \emph{open} otherwise -- for instance, the causal
history of a depicted configuration in a painting, which is compatible
with an unbounded and non-enumerable set of candidate causes.
\end{definition}

\begin{proposition}[Closed Hidden-State Sufficiency]
\label{prop:closed-sufficiency}
In a domain with a finite hidden-state set $H$, the augmented rendering
$r' = (r, H)$ is sufficient for any task $\tau$ whose correctness set
$\mathcal{C}_\tau(e)$ is determined by distinctions drawn from $H$,
provided all task-relevant elements of $H$ are included.
\end{proposition}

\begin{proofsketch}
Since $H$ is finite, exhaustive enumeration eliminates ambiguity about
which hidden variables are missing: every continuation's membership in
$\mathcal{C}_\tau(e)$ becomes computable from $r'$ once the enumerated
variables are fixed, so $\admis{\tau}{r'} \subseteq \mathcal{C}_\tau(e)$
as required by \cref{def:sufficiency}. Chess serves as a witness
(Chapter~\ref{ch:chess}): a board position augmented with castling rights,
en passant status, and repetition count is sufficient for determining
every legal continuation of the game, whereas the visible board alone is
not.
\end{proofsketch}

\begin{proposition}[Sufficiency Achievable Only in the Closed Case]
\label{prop:sufficiency-closed-only}
Sufficiency in the sense of \cref{prop:closed-sufficiency} -- achieved by
finite augmentation of the hidden state -- is available only in closed
hidden-state domains.
\end{proposition}

\begin{proofsketch}
In an open hidden-state domain (\cref{def:closed-open}), no finite set $H$
exhausts the candidate causes or hidden variables relevant to a given
task, by definition of openness; hence no finite $H$ makes
\cref{prop:closed-sufficiency}'s hypothesis satisfiable. Painting
(Chapter~\ref{ch:painting}) is the canonical witness to this proposition's
converse case, and there $\mathcal{C}_\tau(e)$ is typically large rather
than singleton, which is exactly what allows sufficiency to be achieved by
other means (\cref{prop:benign-preserves-sufficiency}) even though
augmentation specifically is unavailable.
\end{proofsketch}
